This thesis studies finite-horizon stochastic impulse control problems for jump-diffusion processes and considers two approaches for the problem's solution. First, a review of the dynamic programming principle is provided. Second, the problem is analysed through its parabolic quasi‑variational inequality (QVI) formulation. In general, solutions to QVIs are non-smooth, and classical solution may not exist. Therefore, a viscosity solution approach is studied and existence as well as uniqueness results are studied. The QVI is then solved numerically using a finite‑difference scheme.

The impulse control problem and deep reinforcement learning (DRL) are then connected by means of dynamic programming. Two simulation-based DRL algorithms are of particular interest in this thesis, namely the value‑based Double Deep Q‑Network (DDQN) and the policy‑based Proximal Policy Optimisation algorithm augmented with self‑imitation learning (PPO‑SIL).

These algorithms are applied to two financial problems. The first considers optimal control of a foreign exchange (FEX) rate with transaction costs, where the DDQN and PPO-SIL algorithms are compared against the finite-difference scheme approximation. The second addresses the dynamic portfolio allocation problem with transation costs, where PPO-SIL is used to trade a 20-asset equity portfolio over a one-year horizon. Asset prices are simulated using correlated multidimensional Variance–Gamma processes with common subordinators and calibrated to European call option prices. In this setting, the learned trading strategy achieves average cumulative returns of approximately 10\% and an average annualised Sharpe ratio of 3.26.